\documentclass[11pt, a4paper]{article}

\usepackage[utf8]{inputenc}
\usepackage[T1]{fontenc}
\usepackage{amsmath}
\usepackage{amssymb}
\usepackage{enumitem}
\usepackage{esvect}
\usepackage{siunitx}
\sisetup{
    locale = DE,
    inter-unit-product = \cdot,
    per-mode = symbol-or-fraction
}

\makeatletter
\def\input@path{{./}{../../}}
\makeatother

\input{defs}

\begin{document}

\section*{Nicht verwendete, aber grundsätzlich gute Beispiele}

\subsection*{Differentialgleichung durch Integration II}

Für die Anzahl $N(t)$ nicht zerfallener Kerne gelte die Differentialgleichung
\[
    \frac{\dd N}{\dd t} = -\lambda N
\]
mit $\lambda > 0$ und der Anfangsbedingung $N(0)=N_0$.
\begin{enumerate}[label=(\alph*)]
    \item Lösen Sie die Differentialgleichung durch Trennung der Variablen.
    \item Bestimmen Sie die Zeit $t_{1/2}$, bei der $N(t_{1/2}) = N_0/2$ gilt.
\end{enumerate}

\subsection*{Freier Fall und Energieerhaltung}

Eine Kugel wird aus der Höhe $h = \SI{5}{\meter}$ aus der Ruhe fallengelassen. Der Luftwiderstand wird vernachlässigt.
\begin{enumerate}[label=(\alph*)]
    \item Wie groß ist ihre potentielle Energie vor dem Loslassen?
    \item Welche Geschwindigkeit hat die Kugel unmittelbar vor dem Aufprall?
\end{enumerate}

\subsection*{Vorzeichen der Arbeit}

Ein Schlitten bewegt sich auf einer \textbf{waagrechten} Strecke nach rechts. Auf ihn wirken eine Zugkraft nach rechts, eine Reibungskraft nach links, die Gewichtskraft nach unten und die Normalkraft nach oben.
\begin{enumerate}[label=(\alph*)]
    \item Welche Kraft verrichtet positive Arbeit?
    \item Welche Kraft verrichtet negative Arbeit?
    \item Welche Kräfte verrichten keine Arbeit?
\end{enumerate}
Begründen Sie Ihre Antworten jeweils kurz.

\subsection*{Beispiel 8 — Arbeit entlang einer Raumkurve I}

Auf ein Teilchen wirke die ortsabhängige Kraft
\[
    \ivec{F}(x,y) = \inlrowTwo{2x}{3}\,\si{\newton}
\]
und das Teilchen bewege sich entlang der Bahn
\[
    \ivec{r}(t) = \inlrowTwo{t}{2t}\,\si{\meter}
    \qquad \text{für } t \in [0,2].
\]
Berechnen Sie die verrichtete Arbeit mit dem Kurvenintegral
\[
    W = \int \ivec{F}\cdot \dd \ivec{r}.
\]
\textit{Hinweis:} Setzen Sie zuerst die Bahn in das Kraftfeld ein und bestimmen Sie anschließend $\dd \ivec{r}/\dd t$.

\subsection*{Beispiel 9 — Arbeit entlang einer Raumkurve II}

Eine Kraft sei gegeben durch
\[
    \ivec{F}(x,y) = \icolTwo{x+y}{2x-y}\,\si{\newton}.
\]
Ein Körper bewege sich auf der Kreisbahn
\[
    \ivec{r}(\varphi) = R\,\icolTwo{\cos\varphi}{\sin\varphi}
    \qquad \text{mit } R = \SI{2}{\meter},\; \varphi \in \left[0,\frac{\pi}{2}\right].
\]
Berechnen Sie die Arbeit
\[
    W = \int \ivec{F}\cdot \dd \ivec{r}
\]
auf diesem Viertelkreis.

\end{document}
